\chapter{Color \& styling}\label{ch:color}
\begin{goals}
\textbf{In this chapter} you color a layer by any of its columns, set the range of a ramp, turn a grade into classes with breaks you place on its histogram, tidy a category legend, decide how blanks look, and carry a finished look to another layer, model or project.
\end{goals}

Color is how a model answers a question at a glance: where is the grade, which rock is it in, what did I flag. In micro a color is always \emph{a column seen through a style}. The column can be anything the layer carries — a source column such as \col{FE}, a calculated one (\cref{ch:filter}), a painted one (\cref{ch:picking}). The style says how its values become colors: a continuous ramp, a set of classes, a legend of categories, or a list of rules.

Nothing in this chapter changes your data. A style is display state: it is saved with the project, and it is what the legend of a figure draws (\cref{ch:figures}). A scene (\cref{sec:sections-scenes}) remembers which column each layer is colored by, and because a style is stored against its column, that column's ramp and classes come back with it.

\section{Where color lives}\label{sec:color-where}
Color has two controls, and they show the same state.

The \ui{color} dropdown, on the second toolbar row, is the quick one. It drives the \emph{active} layer — the one highlighted in the layers panel — and lists what that layer can be colored by: elevation, each numeric column, the category column, \ui{rules}, and every derived column, marked by kind (\code{ƒ} calculated, \code{✎} painted, \code{\%} a fraction-inside column). A calculated column is offered only when it is numeric: a text \code{ƒ} filters and exports, but cannot color (\cref{sec:filter-colorby}). Beside it sit the \ui{min} and \ui{max} boxes that clip the color scale.

\menu{Properties…>style} is the full one. Open it from a layer's right-click menu, from \menu{Data>Layer properties…}, or with \key{⇧P}.\sidekey{⇧P}{properties of the active layer} It asks two things first: a \ui{mode} and a \ui{column}.

\begin{center}
\begin{tabularx}{\linewidth}{@{}lL@{}}\toprule
mode & what it does \\ \midrule
\ui{continuous} & a numeric column (or elevation) through a ramp, optionally cut into classes \\
\ui{categorized} & a text column as a legend of classes, one color each \\
\ui{rule-based} & ordered expressions; the first rule a record matches colors it \\
\ui{painted} & a painted column, in the colors of its palette (offered once the layer has one) \\
\ui{solid} & one color, nothing to map \\ \bottomrule
\end{tabularx}
\end{center}

That list is for block models and drillholes. A point cloud offers \ui{continuous}, \ui{categorized} (its classification) and \ui{rgb} (the colors stored in the file).

\begin{note}
Style edits apply as you make them. On a very large layer each edit costs a redraw, so the footer of the properties panel has an \ui{auto} checkbox: untick it, make several edits, and press \ui{apply} to push them together. \ui{revert} returns to the style the panel opened with.
\end{note}

\shot{color-props-style}{The style tab for the demo deposit, colored by \col{FE} with the clip left automatic: mode and column on top, the ramp cut into six classes by five breaks (each row with its count and share), the no-data row, and the histogram that edits them — the waste piled up around 20\,\% Fe, the ore a long, low tail.}

\section{A ramp and its range}\label{sec:color-ramp}
In \ui{continuous} mode, the \ui{ramp} section picks the palette. There are six.

\begin{center}
\begin{tabularx}{\linewidth}{@{}llL@{}}\toprule
ramp & kind & a good use \\ \midrule
\code{viridis} & sequential & the default; grades, elevation \\
\code{magma} & sequential & grades on a light background \\
\code{turbo} & sequential & wide ranges where you want many distinguishable hues \\
\code{greys} & sequential & a quiet backdrop under colored drillholes \\
\code{spectral} & diverging & a quantity with a meaningful middle \\
\code{rdbu} & diverging & differences: white sits exactly in the middle, so a range symmetric about zero centers on it \\ \bottomrule
\end{tabularx}
\end{center}

A ramp stretches from the column's lowest value to its highest unless you set a \term{clip}: type a minimum and a maximum in the \ui{clip} row (or in the \ui{min}/\ui{max} boxes on the toolbar — they are the same boxes). Everything at or below the minimum takes the first color, everything at or above the maximum the last. An empty box means automatic.

Clip when a few outliers eat the ramp. In the demo deposit nine blocks in ten are waste between 12 and 35\,\% Fe, and an automatic ramp from 12 to 68 spends nearly half its colors on them. A clip of 40–65 gives the colors to the ore.

Two details make the clip safe to use:

\begin{itemize}
\item \textbf{The clip belongs to the column.} Set 40–65 on \col{FE}, switch the color to \col{SG}, and \col{SG} gets its own automatic range. Switch back and \col{FE}'s 40–65 returns.
\item \textbf{The legend says it is clipped.} A clipped end is labeled \code{≤ 40} or \code{≥ 65}, so a reader of your figure knows the darkest color is “40 or less”, not “40”.
\end{itemize}

A calculated (\code{ƒ}) or stored column ramps the same way. Its clip is typed in the column's own units, like a source column's, and a block where it is blank draws in a dark no-data gray, never in the ramp's lowest color.

\begin{tryit}
Open the demo project (on the empty screen, \ui{open the demo project}). The deposit arrives colored by \col{FE}.
\begin{enumerate}
\item Click the deposit in the layers panel to make it active, then press \key{⇧P} and choose the \ui{style} tab.
\item In the \ui{clip} row type \code{40} and \code{65}. The waste collapses into one dark color, and the lens shows its hematite core.
\item Change the ramp to \code{magma}, then to \code{turbo}. The clip stays.
\item In the \ui{column} dropdown choose \col{SG}. The clip boxes empty: this is another column. Choose \col{FE} again and 40–65 comes back.
\end{enumerate}
\end{tryit}

\section{Classes: breaks on the histogram}\label{sec:color-classes}
A continuous ramp is good for looking. A report wants \term{classes}: below 20, 20 to 40, 40 to 55, and so on, each a flat color with a name in the legend. In micro, classes are a property of the ramp.

Type the breaks into the \ui{breaks} box, separated by commas — \code{20, 40, 55, 60, 65} — and tick \ui{classed}. Five breaks make six classes, labeled \code{< 20}, \code{20 – 40}, \code{40 – 55}, \code{55 – 60}, \code{60 – 65} and \code{≥ 65}. Each class is a row with a color chip you can set, taken from the ramp to begin with.

With breaks typed but \ui{classed} \emph{unticked}, the ramp stays continuous and bends at the breaks: each interval gets an equal share of the palette. Use it to give a narrow ore range more color without making classes.

\begin{figure}[htbp]\centering
\definecolor{cA}{HTML}{440154}\definecolor{cB}{HTML}{3B528B}\definecolor{cC}{HTML}{21918C}
\definecolor{cD}{HTML}{5EC962}\definecolor{cE}{HTML}{B8DE29}\definecolor{cF}{HTML}{FDE725}
\begin{tikzpicture}[x=1.55mm, y=1mm, font=\sffamily\small]
% value axis 0..70 (% Fe)
% histogram bars (schematic bimodal grade distribution), 5-unit bins
\foreach \binx/\binh/\binc in {0/30/cA, 5/36/cA, 10/24/cA, 15/13/cA, 20/8/cB, 25/6/cB, 30/5/cB, 35/7/cB,
                      40/10/cC, 45/14/cC, 50/19/cC, 55/24/cD, 60/21/cE, 65/11/cF}
  \fill[\binc!70] (\binx+0.35, 0) rectangle (\binx+4.65, \binh);
\draw[rule, line width=0.5pt] (0,0) -- (71,0);
\foreach \tick in {0,10,...,70} { \draw[rule] (\tick,0) -- (\tick,-1.2); \node[text=muted, font=\sffamily\footnotesize, anchor=north] at (\tick,-1.4) {\tick}; }
\node[text=muted, font=\sffamily\footnotesize, anchor=north] at (35,-6) {FE (\%)};
% breaks
\foreach \brk in {20,40,55,60,65} {
  \draw[amber, line width=0.9pt] (\brk,-0.2) -- (\brk,41);
  \fill[amber] (\brk,41) circle[radius=1.6pt];
}
\node[text=amber, anchor=south, font=\sffamily\footnotesize] at (20,42) {20};
\node[text=amber, anchor=south, font=\sffamily\footnotesize] at (40,42) {40};
\node[text=amber, anchor=south, font=\sffamily\footnotesize] at (54.2,42) {55};
\node[text=amber, anchor=south, font=\sffamily\footnotesize] at (60,42) {60};
\node[text=amber, anchor=south, font=\sffamily\footnotesize] at (65.8,42) {65};
\node[text=ink, anchor=west, font=\sffamily\footnotesize, align=left] at (72,34) {breaks —\\drag to move};
% the class strip
\begin{scope}[yshift=-17mm]
\fill[cA] (0,0) rectangle (20,5); \fill[cB] (20,0) rectangle (40,5); \fill[cC] (40,0) rectangle (55,5);
\fill[cD] (55,0) rectangle (60,5); \fill[cE] (60,0) rectangle (65,5); \fill[cF] (65,0) rectangle (71,5);
\draw[rule, line width=0.4pt] (0,0) rectangle (71,5);
\node[text=muted, anchor=north, font=\sffamily\footnotesize] at (10,-0.8) {$<$ 20};
\node[text=muted, anchor=north, font=\sffamily\footnotesize] at (30,-0.8) {20 – 40};
\node[text=muted, anchor=north, font=\sffamily\footnotesize] at (47.5,-0.8) {40 – 55};
\node[text=muted, anchor=north, font=\sffamily\footnotesize] at (57.5,-6) {55 – 60};
\node[text=muted, anchor=north, font=\sffamily\footnotesize] at (62.5,-0.8) {60 – 65};
\node[text=muted, anchor=north, font=\sffamily\footnotesize] at (69,-6) {≥ 65};
\draw[rule] (57.5,-0.6) -- (57.5,-5.6); \draw[rule] (68,-0.6) -- (68.6,-5.6);
\node[text=ink, anchor=west, font=\sffamily\footnotesize, align=left] at (72,2.5) {classes —\\one color each};
\end{scope}
\end{tikzpicture}
\caption{Five breaks cut a grade into six classes (schematic). A class is its interval: it keeps its color wherever the data's range, the clip, or the layer puts it.}\label{fig:color-classes-diagram}
\end{figure}

\subsection{The histogram is the editor}
Typing breaks works on any numeric column. Placing them well needs the distribution, and micro shows it once the column has been \emph{scanned}. If it hasn't, the ramp section offers a button, \ui{scan FE for its histogram + quantiles}; one pass over the table fills it in (\cref{ch:columns} has more on the scan).

Then the histogram under the ramp is live:

\begin{itemize}
\item \textbf{Drag} a break to move it. The 3D view recolors as you drag, and the class rows update their bounds, counts and shares.
\item \textbf{Double-click} to add a break. Both halves of the split class keep its color until you change one.
\item \textbf{Right-click} a break to remove it. The lower class's color stays.
\item \ui{log} switches the count axis to a log scale. A grade distribution buries its ore tail on a linear axis; this brings it up. micro ticks it for you when one bin towers over the rest.
\end{itemize}

Next to \ui{log} are shortcuts that make breaks for you. Type a class count (2 to 12, five to begin with) and click \ui{quantiles} for classes holding equal numbers of records, or \ui{equal} for classes of equal width across the ramp's range; \ui{clear} removes every break. When one value dominates — waste at the lowest grade — the quantiles would collapse onto it, so micro takes them from the values above it instead.

Every break edit is one undo step (\key{⌘Z}, or \key{Ctrl}+\key{Z} on Windows and Linux).

\subsection{A class is its interval}
A class is defined by its bounds, and its color belongs to those bounds. Change the clip, reload the model with a wider range, look at a filtered subset: \code{55 – 60} is still the same green.

It also means two layers can share a legend. Give the drillhole assays the same breaks and colors as the block model (\cref{sec:color-copy} does it in two clicks) and a hole reads against the blocks around it, color for color.

\subsection{Hiding a class}
Each class row has an eye. A hidden class is not drawn, cannot be picked, and leaves the legend. Hiding everything below 40 is a fast way to see the ore body alone.

\begin{caution}
A hidden class is a \emph{display} choice. The blocks are still in the model, and grade–tonnage, stats and exports still count them. To restrict an analysis, use the filter (\cref{ch:filter}) or a selection.
\end{caution}

\begin{tryit}
On the demo deposit, colored by \col{FE}, in the \ui{style} tab (if you skipped the last exercise: open the demo project, click the deposit, press \key{⇧P}):
\begin{enumerate}
\item Empty both \ui{clip} boxes, so the histogram spans every value.
\item In the \ui{breaks} box type \code{20, 40, 55, 60, 65} and tick \ui{classed}. Six class rows appear.
\item If the histogram is empty, click the \ui{scan FE…} button and wait for the pass to finish.
\item Drag the break at 55 slowly to the left. The cut face of the model recolors as you go, and the class rows recount; stop near 50.
\item Double-click the histogram near 30. \code{20 – 40} splits in two, both halves the same color.
\item Click the color chip of the new upper class and choose a different color.
\item Click the eye of \code{< 20}. The poorest waste — 22\,\% of the blocks — disappears from the view and from the legend.
\item Press \key{⌘Z} until the model is back to five breaks, everything visible.
\end{enumerate}
\end{tryit}

\section{Categories and their legend}\label{sec:color-categories}
\ui{categorized} mode colors by a text column: \col{LITO} in the demo, with HEMATITE, ITABIRITE, CANGA and WASTE. The \ui{column} dropdown picks which column classes the layer; any column can, zone codes stored as numbers included. On a model, changing it re-reads the file once.

Below the column is the class list, also reachable as the \ui{classes} tab. Each entry has an eye, a color chip, a name and a count with its share, which micro fills in by reading the layer once. Lists longer than eight entries get a search box, and \ui{all}, \ui{none} and \ui{invert} act on every eye at once.

The legend is yours to tidy:

\begin{itemize}
\item \textbf{Recolor}: click the chip.
\item \textbf{Reorder}: drag the \code{⠿} grip. The order is the order of the figure's legend.
\item \textbf{Rename}: right-click an entry, \menu{Rename…}. The data keeps its value; the legend shows your label.
\item \textbf{Merge}: right-click, \menu{Merge into}, and pick another entry. Merged classes share one color and one eye — ITABIRITE and a banded-iron code under one “Iron Fm”, say. \menu{Ungroup} separates them again.
\item \textbf{Remove}: right-click, \menu{Remove (→ other)}. The values fall into a gray \emph{other} entry; \menu{Split into entries} on \emph{other} brings them back.
\item \ui{reset legend}, offered once you have edited the legend, returns to the raw values.
\end{itemize}

\begin{note}
Renaming and merging here change the \emph{legend}, never the column. A filter still says \code{LITO = "ITABIRITE"}, and an export still writes ITABIRITE. To change the data, derive a new column (\cref{sec:filter-calc}).
\end{note}

As with grade classes, a hidden category is not drawn, not pickable and not in the legend — and still counted by every analysis.

\shot[0.62\linewidth]{color-categories}{The class list for \col{LITO} on the demo deposit: an eye, a color and a count for each lithology — CANGA 36,495, HEMATITE 11,013, ITABIRITE 129,196 — with WASTE (1,609,916 blocks, 90\,\%) hidden.}

\section{Rules}\label{sec:color-rules}
Sometimes neither one grade nor one category says what you mean: “high grade \emph{in hematite}”, “anything flagged, whatever its grade”. \ui{rule-based} mode colors by expressions, written in the same language as the filter (\cref{ch:filter}):

\begin{lstlisting}[language=expr]
FE >= 60 and LITO = "HEMATITE"
FE >= 55
LITO = "CANGA"
\end{lstlisting}

Rules are ordered and \textbf{the first match wins}. A hematite block at 62\,\% takes the first rule's color; one at 57\,\% falls to the second. Whatever matches nothing takes the color of the last row, \ui{everything else}.

Each rule has a name (it is what the legend shows; unnamed, the expression stands in), a color, an eye, and a count. \ui{+ add rule} appends one, and the \code{⠿} grip drags a rule to a new priority. The \code{⋮} button on a row holds \menu{Duplicate} and \menu{Remove}.

Rules are evaluated in one sweep over the table. Colors, names and eyes change instantly. Editing an expression, or adding, removing or reordering a rule, changes what the sweep means, so the \ui{run rules} button shows a dot until the sweep has run again. With \ui{auto} on, it runs on its own. Rules are saved with the project and swept again when it opens.

\begin{tip}
Rules are the right tool when the legend should read as sentences. If what you want is a column you can also filter, tabulate and export, write the same logic once as a calculated column that returns a class code — \code{if(LITO = "HEMATITE", 2, if(FE >= 55, 1, 0))} — and color it with breaks between the codes (\cref{sec:filter-calc}). The column must return a number to color; a text result filters and exports, but is not offered in the \ui{color} list.
\end{tip}

\section{Blanks: the no-data color}\label{sec:color-nodata}
A blank grade is \emph{unestimated}. It is not zero and it is not the lowest class, so micro never draws it as either. Blanks draw in their own color, gray by default.

The \ui{no data} row of the style tab sets that color, or hides blanks with its \ui{hide} checkbox. Once the column has been scanned the row also says how many there are — \code{12,408 blank}, say (the demo deposit reads \code{0 blank}) — and that count is worth a look before you trust a mean. If any are visible, the figure legend lists a \emph{no data} entry.

A value that only \emph{means} blank is a different matter. A \code{-99} written by an estimator is a number until you say otherwise, and it will drag the ramp and every statistic down. Declare it once: in \menu{Properties…>columns}, right-click the column and choose \menu{Detection limits \& blanks>Values that mean blank…}. That works for CSV and drillhole tables; Datamine and Parquet files carry real nulls of their own. \Cref{ch:columns} covers blanks and detection limits in full.

\section{Opacity, edges, and how things are drawn}\label{sec:color-draw}
A few display settings are not about color, yet they decide how well it reads.

\begin{description}
\item[Opacity] Right-click a layer, \menu{Opacity}, and pick 100, 75, 50 or 25\,\%; the style tab has a slider from 5 to 100\,\%. It works on block models, drillholes and meshes, and a group's menu sets it for all its members. micro draws transparency as a fine dither, so a see-through shell never hides what is behind it in the wrong order.
\item[Block edges] \key{⇧B}, or \menu{View>Block edges}, outlines every block — and every cut block on a section. It is how you check sub-blocking and lattice alignment. A layer can override the global switch: right-click, \menu{Block edges}, and choose \ui{auto} (follow \ui{View}), \ui{on} or \ui{off}.
\item[Boxes or points] \menu{View>Blocks as boxes} draws true block faces; \menu{View>Blocks as points} draws each block as a splat, which is lighter on a slow machine.
\item[Point size] A layer drawn as points — a cloud, a model opened as points — has its own \menu{Point size} (×1, ×2, ×3, ×5, ×8), a multiple of the global \ui{pt px} slider.
\item[Stick radius] Drillholes draw as sticks; \menu{Stick radius…} or the \ui{stick radius (m)} row sets how fat (\cref{ch:drillholes}).
\item[Tint] A mesh has no column to color by here; \menu{Tint…} gives it one color from a short geological palette (Grey, Sand, Rust, Green, Blue, Violet, Brown) or a custom hex value (\cref{ch:meshes}).
\end{description}

The \ui{EDL} checkbox and its \ui{strength} slider control the shading that makes a point cloud read as a surface, and \ui{budget} sets how many elements are drawn per frame: lower is snappier while you move, and the picture converges when you stop.

\section{Reusing a look}\label{sec:color-copy}
Once the classes, colors and no-data choices are right on one layer, you should not have to build them again. Right-click the layer and open \menu{Style}:

\begin{center}
\begin{tabularx}{\linewidth}{@{}lL@{}}\toprule
\menu{Copy style} & takes the layer's whole look. The copy also goes to the system clipboard, so it survives into another project or another tab. \\
\menu{Paste style} & applies it to the layer you right-clicked. \\
\menu{Save style…} & writes it as a \file{.microstyle.json} file. \\
\menu{Load style…} & applies such a file. \\ \bottomrule
\end{tabularx}
\end{center}

Everything is matched \textbf{by column name}. The \col{FE} classes land on \col{FE} wherever \col{FE} sits in the other file; a \col{LITO} legend lands on a column called \col{LITO}. The status line then says what applied and what had no match, so a silent half-paste cannot happen.

Three uses come up again and again: the drillholes in the model's grade classes; next month's model in this month's legend; a company standard kept as one \file{.microstyle.json} that every project loads. \menu{Change file…}, which points a layer at a new file, keeps the look by column name the same way.

\begin{tryit}
With the demo deposit classed as in the last exercise (if you skipped it, type \code{20, 40, 55, 60, 65} in the \ui{breaks} box and tick \ui{classed}):
\begin{enumerate}
\item Right-click the deposit, \menu{Style>Copy style}.
\item In the layers panel, open the drillhole set and right-click its assay layer. Choose \menu{Style>Paste style}.
\item Read the status line: \emph{style from deposit.csv → holes · demo\_assay.csv: FE, colour by FE}. The sticks now wear the model's classes, and where a hole crosses the cut face the colors agree — or show you where they don't. (In the demo they agree: the holes were sampled from the same rock as the blocks.)
\end{enumerate}
\end{tryit}

\section{The legend, the background and the theme}\label{sec:color-legend}
The legend draws what the model draws. Classes appear as swatches in their colors; a continuous ramp as its bar, with clipped ends marked; categories under their labels, hidden ones left out; rules under their names. It is titled with the column's name, and two legends on the same column say which layer each belongs to.

Turn it on in \menu{View>Decorations \& figure…}, where each colored layer has its own checkbox under \ui{legends}, next to the scale bar, north arrow and title. \Cref{ch:figures} covers that panel. You never draw a legend by hand, so it cannot disagree with the picture.

The same panel sets the \ui{background}: \ui{match the theme}, \ui{basalt}, \ui{black}, \ui{white}, \ui{paper}, or a custom color. A white or paper background is what a report page wants, and the decorations switch to dark ink on it. Check your ramp against the background you will publish on: \code{viridis} ends in a yellow that fades on white, where \code{magma} or \code{turbo} hold.

The interface itself has a light and a dark \term{theme}, under \menu{View>Theme}, which can also follow the system or be a custom theme you edit. The theme is about the panels and menus; with the background on \ui{match the theme}, the 3D view follows it.

\begin{why}
The clip, the classes, the colors and the no-data choice are all stored against the \emph{column}, not against the layer's current state. That is why switching the color dropdown back and forth loses nothing, why a pasted style can find its target by name, and why a figure made in March and one made in June use the same green for 55–60\,\% Fe without anyone remembering to make them match.
\end{why}
